% Options for packages loaded elsewhere
\PassOptionsToPackage{unicode}{hyperref}
\PassOptionsToPackage{hyphens}{url}
\documentclass[
]{article}
\usepackage{xcolor}
\usepackage{amsmath,amssymb}
\setcounter{secnumdepth}{-\maxdimen} % remove section numbering
\usepackage{iftex}
\ifPDFTeX
  \usepackage[T1]{fontenc}
  \usepackage[utf8]{inputenc}
  \usepackage{textcomp} % provide euro and other symbols
\else % if luatex or xetex
  \usepackage{unicode-math} % this also loads fontspec
  \defaultfontfeatures{Scale=MatchLowercase}
  \defaultfontfeatures[\rmfamily]{Ligatures=TeX,Scale=1}
\fi
\usepackage{lmodern}
\ifPDFTeX\else
  % xetex/luatex font selection
\fi
% Use upquote if available, for straight quotes in verbatim environments
\IfFileExists{upquote.sty}{\usepackage{upquote}}{}
\IfFileExists{microtype.sty}{% use microtype if available
  \usepackage[]{microtype}
  \UseMicrotypeSet[protrusion]{basicmath} % disable protrusion for tt fonts
}{}
\makeatletter
\@ifundefined{KOMAClassName}{% if non-KOMA class
  \IfFileExists{parskip.sty}{%
    \usepackage{parskip}
  }{% else
    \setlength{\parindent}{0pt}
    \setlength{\parskip}{6pt plus 2pt minus 1pt}}
}{% if KOMA class
  \KOMAoptions{parskip=half}}
\makeatother
\setlength{\emergencystretch}{3em} % prevent overfull lines
\providecommand{\tightlist}{%
  \setlength{\itemsep}{0pt}\setlength{\parskip}{0pt}}
\usepackage[]{natbib}
\bibliographystyle{plainnat}
\newcommand{\bigsqcap}{\mathop{\Large\sqcap}\limits}
\usepackage{stmaryrd}
\usepackage{bookmark}
\IfFileExists{xurl.sty}{\usepackage{xurl}}{} % add URL line breaks if available
\urlstyle{same}
\hypersetup{
  pdftitle={Constructing quotient inductive types from ordinals},
  hidelinks,
  pdfcreator={LaTeX via pandoc}}

\title{Constructing quotient inductive types from ordinals}
\author{\shortstack{Thorsten Altenkirch and Christina O'Donnell\tabularnewline University of Nottingham}}
\date{}

\begin{document}
\maketitle

\subsection{Abstract}\label{abstract}

Quotient inductive types extend inductive types with equations between
constructors, supporting definitions such as multisets, infinitary
unordered trees, extensional ordinals, and internal syntax. Although
such types are readily specified, their construction from more
elementary type formers is subtle: existing constructions of general
infinitary QW-types use the Weakly Initial Set of Covers axiom to
establish a cocontinuity property.

We construct a class of infinitary QW-types, including mobiles, finite
multisets, and the cumulative hierarchy of extensional sets, from
W-types and quotients relative to an extensional plump ordinal basis,
without assuming WISC. Our construction applies to depth-preserving
signatures, whose equations preserve an ordinal-valued rank on terms.
This rank descends through the quotient and yields a section of the
canonical map associated with the colimit of the inflationary iteration,
thereby establishing the required cocontinuity.

More generally, we isolate the existence of such a section as the
semantic criterion underlying the construction. Depth preservation
provides a syntactic sufficient condition, while normalising and
finitary presentations give further instances of the same principle.

\section{1. Introduction}\label{introduction}

Inductive types provide a fundamental method for defining mathematical
objects in dependent type theory. Their constructors generate the
elements of a type freely, giving rise to recursion and induction
principles. Many natural definitions, however, are not free: their
constructors must additionally satisfy equations. Quotient inductive
types (QITs) extend inductive types by allowing such equations to be
specified alongside the constructors
\citep{altenkirch2016-type-theory-qit}. Examples include finite
multisets, permutation-invariant infinitary trees, extensional ordinals,
cumulative hierarchies of extensional sets, and internal representations
of the syntax of type theory. They can also be understood as the
set-truncated fragment of higher inductive types.

Although QITs are readily specified, it is not generally clear whether
they can be constructed from more elementary type formers. In classical
set theory, an equational inductive type can often be obtained by first
constructing the corresponding type of raw well-founded trees and then
quotienting by the equations. Constructively, this procedure need not
produce the required initial algebra. The difficulty is particularly
apparent for infinitary signatures: defining the algebra structure on
the quotient can require the simultaneous choice of representatives for
an entire family of quotient elements.

Fiore, Pitts and Steenkamp isolate a broad algebraic class of QITs
called QW-types, and its indexed generalisation, QWI-types. A QW-type is
the initial algebra for an equational theory whose operations may have
infinitary arities. They show that QW-types can be constructed from
W-types and quotients in the presence of the Weakly Initial Set of
Covers axiom (WISC) \citep{fiore2022-quotient-inductive}. W-types
originate in Martin-Löf's constructive type theory
\citep{martin-lof1984-mltt}. WISC appears as a collection principle in
site-theoretic and realizability settings
\citep{roberts2012-internal-categories, streicher2005-wisc}. Their
construction uses \emph{inflationary iteration}
\citep{pitts2021-inflationary-iteration}: rather than forming a single
quotient of the type of raw terms, it constructs a diagram of successive
approximations in which applications of constructors and equations are
introduced stage by stage. The required QW-type is then obtained as the
colimit of this diagram.

The central issue in this construction is a cocontinuity property. Let
\(S\) be the polynomial functor determined by the operations of a
QW-signature, represented by a container in the sense of Abbott,
Altenkirch and Ghani \citep{abbott2005-container-derivative}, and let
\(D\) be the diagram of approximations. There is a canonical map

\[
\kappa_D :
\mathrm{colim}_\alpha S (D_\alpha)
\longrightarrow
S \left(\mathrm{colim}_\alpha D_\alpha\right).
\]

To equip the colimit with the required algebra structure, one must lift
a family of elements of the colimit to some common stage of the diagram.
For infinitary operations, there need not be an evident stage containing
representatives of every member of the family. WISC is used by Fiore,
Pitts and Steenkamp to construct a sufficiently rich notion of size and
to prove the corresponding cocontinuity result.

The obstruction can be seen concretely in the QW-type of \emph{mobiles}.
Informally, a mobile is a well-founded countably branching tree whose
children are unordered. It is generated by an operation
\(\mathsf{node} : (\mathbb{N} \to M) \to M\) subject to equations
identifying \(\mathsf{node} (f)\) and \(\mathsf{node} (f \circ \pi)\)
for every permutation \(\pi\) of \(\mathbb{N}\). To place
\(\mathsf{node} (f)\) at a stage of an inflationary construction, we
must find a stage containing all the children \(f (n)\). But elements of
the colimit are represented by elements at possibly different stages,
and quotienting removes the particular representatives from which such
stage information might be recovered. Moreover, two enumerations
representing the same unordered family need not induce definitionally
equal sizes when sizes are themselves inductively generated. Thus the
desired stage assignment must be invariant under the equations of the
QW-type.

In this paper, we obtain such an invariant stage assignment from an
\emph{extensional plump ordinal basis}. Plump ordinals provide
sufficiently large upper bounds for families of smaller ordinals, while
extensionality ensures that an ordinal is determined by its strict
predecessors. Consequently, the rank assigned to an infinitary
constructor can depend extensionally on the ranks of its arguments,
rather than on a particular enumeration or representation of those
arguments.

Swan's W-types with reductions provide a related extension of polynomial
initial algebras and connect them to the small object argument
\citep{swan2018-w-types-reductions}. His later ZF construction of
image-preserving higher inductive types similarly obtains
representative-independent stages without choice
\citep{swan2022-hits-in-zf}.

We identify a syntactic class of equational theories for which this rank
is compatible with the quotient. We call these theories
\emph{depth-preserving}. Roughly, an equation is depth-preserving when
its two sides assign the same construction depth to every substitution
instance. The recursively defined rank on raw terms is therefore
invariant under the generating equations and descends to the quotient.

Our first main result may be stated informally as follows.

\begin{quote}
\textbf{Theorem.} Let \(\Sigma\) be a depth-preserving QW-signature, and
suppose that there is an extensional plump ordinal basis closed under
the arities of \(\Sigma\). Then the QW-type presented by \(\Sigma\) can
be constructed from W-types and quotients by inflationary iteration,
without use of WISC.
\end{quote}

The descended rank assigns to every element of the quotient a canonical
stage at which it can be represented. Applying this rank to an element
of \(S \left(\mathrm{colim}_\alpha D_\alpha\right)\) produces a stage containing a
lift of all its arguments. This defines a map

\[
\sigma_D :
S \left(\mathrm{colim}_\alpha D_\alpha\right)
\longrightarrow
\mathrm{colim}_\alpha S (D_\alpha)
\] such that \(\kappa_D \circ \sigma_D = \mathsf{id}\).

Thus \(\sigma_D\) is a section of the canonical map \(\kappa_D\). This
section supplies precisely the lifting operation required in the
cocontinuity argument and hence allows the inflationary construction to
proceed. The class of depth-preserving signatures includes mobiles,
standard presentations of finite multisets, and the cumulative hierarchy
of extensional sets.

The semantic criterion is the existence of a coherent section of
\(\kappa_D\). Depth preservation supplies one syntactic route to such a
section; normal forms and finitary branching supply others. This
separates the stage information required by inflationary iteration from
the particular mechanism used to obtain it.

The remainder of the paper is organised as follows. We first recall
QW-signatures, their algebras, and the inflationary construction of
Fiore, Pitts and Steenkamp. We then introduce extensional plump ordinals
and explain their role as representative-independent construction
stages. Next, we define depth-preserving signatures and prove that their
structural ranks descend through the generated equations. We use these
ranks to construct a section of the canonical colimit map and derive the
existence of the corresponding QW-types. Finally, we discuss alternative
sufficient conditions, including normalisation and finitary branching,
and the prospects for a broader syntactic characterisation of
sectionability.

\section{2. Motivating examples}\label{motivating-examples}

Three examples delimit the intended class. Their depth-preservation
proofs are given uniformly in Section 6.3.

\subsection{2.1 Finite multisets}\label{finite-multisets}

For a type \(A\), finite multisets are lists quotiented by permutations,
generated by \(\mathsf{nil}:M (A)\) and
\(\mathsf{cons}:A\to M (A)\to M (A)\). Adjacent transpositions impose
\(\mathsf{cons} (a,\mathsf{cons} (b,x)) = \mathsf{cons} (b,\mathsf{cons} (a,x))\).
Both sides place \(x\) beneath two constructor layers, so list length
descends to the quotient. This finitary example does not need ordinal
lifting, but it isolates the principle that equations must preserve
enough stage information.

\subsection{2.2 Mobiles}\label{mobiles}

A mobile is a well-founded countably branching tree with unordered
children. It has a constructor \(\mathsf{node}:(\mathbb{N}\to M)\to M\)
and equations \(\mathsf{node} (f)=\mathsf{node} (f\circ\pi)\) for
permutations \(\pi:\mathbb{N}\cong\mathbb{N}\); an image-based
presentation identifies any enumerations with the same image.
Constructing a node requires a common stage for all children, but the
colimit supplies only pointwise representatives. Moreover, the common
stage must be independent of their enumeration. An extensional plump
ordinal resolves both problems: plumpness bounds the family of child
ranks, and extensionality identifies bounds generated from the same
image.

\subsection{2.3 The cumulative hierarchy of
sets}\label{the-cumulative-hierarchy-of-sets}

The cumulative hierarchy \(V\) of the HoTT Book is generated by
\(\mathsf{set}:(A:\mathcal{U})\to (A\to V)\to V\) and identifies
\(\mathsf{set} (A,f)\) with \(\mathsf{set} (B,g)\) when \(f\) and \(g\)
have the same image \citep[sec.~10.5]{hott2013-book}. Since every
enumerated element occurs immediately below the outer constructor, image
equality preserves the collection of predecessor ranks. An ordinal basis
closed under the types in \(\mathcal{U}\) then gives a
representative-independent stage even for infinitary \(A\). This example
shows that the method applies to image equations, not only permutations;
see \citep{eleftheriadis2021-cumulative-hierarchy} for the set-theoretic
properties of \(V\).

\subsection{2.4 Other examples and
boundaries}\label{other-examples-and-boundaries}

Reordering, duplication, and extensional reindexing at a fixed
constructor level are typical depth-preserving operations. The condition
is sufficient, not necessary: finitary signatures lift finite families
directly, while normalising signatures may rank canonical normal forms
even when unit or associativity changes raw depth. Section 8 records
these established alternatives.

\section{3. QW-types and their construction
problem}\label{qw-types-and-their-construction-problem}

We now recall the algebraic presentation of QW-types and explain why
their construction from W-types and quotients is non-trivial. We work in
a universe \(\mathcal{U}\) of sets equipped with dependent sums,
dependent products, and W-types \citep{martin-lof1984-mltt}, together
with effective quotients. The additional extensionality and
unique-choice assumptions used by the main theorem are stated in Section
7.

\subsection{3.1 Signatures and
equations}\label{signatures-and-equations}

A polynomial signature is a pair \(\Sigma = (A,B)\) consisting of a type
\(A : \mathcal{U}\) of operation symbols and a family
\(B : A \to \mathcal{U}\) assigning an arity to each operation. An
element \(a : A\) represents an operation whose arguments are indexed by
\(B (a)\). No finiteness assumption is made on \(B (a)\).

The signature \(\Sigma\) determines a polynomial endofunctor
\(S_\Sigma : \mathcal{U} \to \mathcal{U}\), represented by the container
\((A,B)\) \citep{abbott2005-container-derivative}, and defined by
\(S_\Sigma (X) = (a:A)\times (B (a) \to X).\)

For a function \(f : X \to Y\), its action is given by
\(S_\Sigma (f) (a,b) = (a,f \circ b).\)

An \(S_\Sigma\)-algebra is a type \(X : \mathcal{U}\) equipped with a
structure map \(\alpha : S_\Sigma (X) \to X.\)

Equivalently, \(\alpha\) assigns to each \(a : A\) an operation
\(\alpha_a : (B (a) \to X) \to X.\)

A morphism from an algebra \((X,\alpha)\) to an algebra \((Y,\beta)\) is
a function \(h : X \to Y\) satisfying
\(h (\alpha (a,b)) = \beta (a,h \circ b)\) for every \(a : A\) and
\(b : B (a) \to X\).

The W-type \(W_\Sigma\) associated with \(\Sigma\) is the initial
\(S_\Sigma\)-algebra. Its elements may be regarded as closed
well-founded trees whose nodes are labelled by operations \(a : A\) and
whose immediate subtrees are indexed by \(B (a)\).

To describe equations, we also require terms containing variables. For
each type \(X : \mathcal{U}\), let \(T_\Sigma (X)\) be the free
\(S_\Sigma\)-algebra generated by \(X\). It is inductively generated by
constructors \(\eta : X \to T_\Sigma (X)\) and
\(\sigma : S_\Sigma (T_\Sigma (X)) \to T_\Sigma (X).\)

The constructor \(\eta\) inserts a variable, while \(\sigma\) applies an
operation of the signature to a family of terms. The assignment
\(T_\Sigma\) is the free monad generated by \(S_\Sigma\).

Given an \(S_\Sigma\)-algebra \((Y,\alpha)\) and a substitution
\(\rho : X \to Y\), every term \(t : T_\Sigma (X)\) has an
interpretation \(\llbracket t \rrbracket_{\alpha,\rho} : Y\). This
interpretation is defined recursively by
\(\llbracket \eta (x) \rrbracket_{\alpha,\rho} = \rho (x)\) and

\[
\llbracket \sigma (a,b) \rrbracket_{\alpha,\rho}
=
\alpha \left(a,\lambda i.\llbracket b (i) \rrbracket_{\alpha,\rho}\right).
\]

We also write this interpretation as \(t \mathbin{\gg=} \rho\).

A system of equations over \(\Sigma\) consists of a type
\(E : \mathcal{U}\) of equation names, a family
\(V : E \to \mathcal{U}\) of variable contexts, and two families of
terms \(l,r : (e:E)\to T_\Sigma (V (e)).\)

For each \(e : E\), the terms \(l (e)\) and \(r (e)\) are the two sides
of the corresponding equation.

An \(S_\Sigma\)-algebra \((X,\alpha)\) satisfies the equation system if,
for every equation \(e : E\) and every substitution
\(\rho : V (e) \to X\), we have

\[
\llbracket l (e) \rrbracket_{\alpha,\rho}
=
\llbracket r (e) \rrbracket_{\alpha,\rho}.
\]

We write this property as \(\mathsf{Sat}_{\Sigma,E} (X,\alpha).\)

A QW-algebra for the presentation \((\Sigma,E,V,l,r)\) is an
\(S_\Sigma\)-algebra together with a proof that it satisfies all the
specified equations.

QW-types constitute a set-level fragment of higher inductive types. A
higher inductive type may have constructors producing points, paths, and
higher-dimensional paths. A QIT restricts attention to set-truncated
types and equations between their point constructors. QW-types are the
algebraic class of QITs presented by possibly infinitary polynomial
operations and unconditional equations between terms. The subsequent
construction uses only this algebraic presentation and does not
otherwise depend on the higher-dimensional interpretation.

\subsection{3.2 QW-types as initial
algebras}\label{qw-types-as-initial-algebras}

A QW-type for an equational presentation \((\Sigma,E,V,l,r)\) is an
initial QW-algebra. It consists of a type \(Q\), an algebra structure
\(\mathsf{intro} : S_\Sigma (Q) \to Q,\) and a proof
\(\mathsf{equate} : \mathsf{Sat}_{\Sigma,E} (Q,\mathsf{intro}).\)

Initiality states that for every QW-algebra \((X,\alpha,p)\) there is a
unique algebra morphism \(\mathsf{rec}_{X,\alpha,p} : Q \to X.\)

The morphism preserves the constructors:

\[
\mathsf{rec}_{X,\alpha,p} (\mathsf{intro} (a,b))
=
\alpha (a,\mathsf{rec}_{X,\alpha,p} \circ b).
\]

Moreover, if \(h : Q \to X\) is any other function satisfying the same
homomorphism equation, then \(h = \mathsf{rec}_{X,\alpha,p}.\)

This initiality property gives the non-dependent recursion principle of
the QW-type. Under the usual extensionality assumptions, it also yields
the expected dependent elimination principle for the associated quotient
inductive type \citep{fiore2022-quotient-inductive}.

Classically, one might attempt to construct \(Q\) by quotienting the
W-type \(W_\Sigma\) by the smallest congruence containing all
substitution instances of the equations. Let
\(q : W_\Sigma \to W_\Sigma/{\sim}\) denote the resulting quotient map.
To make the quotient into an \(S_\Sigma\)-algebra, we would need a
constructor

\[
\overline{\mathsf{sup}}
:
S_\Sigma (W_\Sigma/{\sim})
\to
W_\Sigma/{\sim}.
\]

An input to this constructor consists of an operation \(a : A\) and a
family \(b : B (a) \to W_\Sigma/{\sim}.\)

The constructor of \(W_\Sigma\) acts only on families of representatives
\(\widetilde{b} : B (a) \to W_\Sigma.\)

We would therefore like to choose \(\widetilde{b}\) satisfying
\(q \circ \widetilde{b} = b\) and define
\(\overline{\mathsf{sup}} (a,b) = q (\mathsf{sup} (a,\widetilde{b})).\)

The congruence ensures that the result is independent of the chosen lift
once such a lift has been obtained. The problem is that the pointwise
surjectivity of \(q\) does not construct a lift of the entire family
\(b\). Such a lift requires a choice principle for families indexed by
\(B (a)\).

For finitary arities, this obstruction can often be resolved by finite
choice. For an infinitary arity \(B (a)\), however, the required lifting
principle may be a genuine additional assumption. In the case of
mobiles, for example, lifting a family
\(\mathbb{N} \to W_\Sigma/{\sim}\) through \(q\) has the form of
countable choice
\citep{altenkirch2016-type-theory-qit, fiore2022-quotient-inductive}.

Categorically, the problem is that the polynomial functor \(S_\Sigma\)
need not preserve quotient maps constructively. The quotient map \(q\)
induces
\(S_\Sigma (q) : S_\Sigma (W_\Sigma) \to S_\Sigma (W_\Sigma/{\sim}),\)
but the surjectivity of \(q\) does not imply the surjectivity of
\(S_\Sigma (q)\) without an appropriate choice principle. Consequently,
the constructor on \(W_\Sigma\) cannot in general be transported
directly across the quotient.

Inflationary iteration avoids asking for all representatives at once.
Instead, it interleaves applications of the polynomial functor with
quotients and forms the candidate QW-type as a colimit of successive
approximations. This replaces the lifting problem for a single quotient
by the problem of showing that \(S_\Sigma\) interacts appropriately with
the resulting colimit. As Section 4.3 explains, this is the cocontinuity
problem addressed by both the WISC-based construction and our
ordinal-rank construction.

\subsection{3.3 Choice strength}\label{choice-strength}

The difficulty above is not merely an artefact of a particular
construction. General infinitary equational theories can have genuine
choice-theoretic strength.

Blass constructed an infinitary equational theory whose free algebras
cannot be proved to exist in ZF set theory without the Axiom of Choice
\citep{blass1983-free-algebra-coequalizer}. Lumsdaine and Shulman
reformulated this theory as a higher inductive type and showed that the
corresponding general existence claim is likewise not provable in ZF
without choice \citep[sec.~9]{lumsdaine2018-higher-inductive-types}.
Thus sufficiently general schemas of infinitary higher or quotient
inductive types are strictly stronger than the unconditional combination
of ordinary inductive types and quotients in a choice-free setting.

This result does not imply that every infinitary QW-type requires
choice. The obstruction depends on the interaction between the
infinitary operations and their equations. Mobiles and the cumulative
hierarchy, for example, have infinitary constructors, but their
equations preserve structural depth. Their elements therefore carry
invariant rank information that is absent from arbitrary infinitary
presentations.

Fiore, Pitts and Steenkamp show that WISC suffices to construct all
QW-types, and more generally QWI-types, covered by their schema
\citep{fiore2022-quotient-inductive}. WISC is substantially weaker than
the full Axiom of Choice and is valid in many constructive models.
Nevertheless, the Blass example shows that some additional restriction
or set-theoretic principle is unavoidable for a completely general
construction.

Our aim is therefore not to construct every infinitary QW-type without
choice. Instead, we identify presentations for which the lifting
required by inflationary iteration can be obtained from intrinsic
structural information. The semantic condition is the existence of a
section of the canonical polynomial-colimit map. Depth preservation
supplies a syntactic condition under which such a section can be
constructed from an extensional plump ordinal basis.

\section{4. Inflationary iteration}\label{inflationary-iteration}

Fiore, Pitts and Steenkamp construct QW-types by an inflationary
iteration in which the formation of terms is interleaved with
quotienting by equations \citep{fiore2022-quotient-inductive,
pitts2021-inflationary-iteration}. This
avoids the direct attempt to lift an entire infinitary family through a
single quotient map. The construction nevertheless requires a
sufficiently strong cocontinuity property, which is the point at which
their argument uses WISC.

Throughout this section, fix a polynomial signature \(\Sigma=(A,B)\) and
an equation system \((E,V,l,r)\) as in Section 3. We write \(T_\Sigma\)
for the free monad generated by the polynomial functor \(S_\Sigma\).

\subsection{4.1 Size-indexed
approximants}\label{size-indexed-approximants}

Let \(\mathsf{Size}\) be a type equipped with a well-founded relation
\(<\). Let \(D\) denote the diagram of approximants. For
\(\alpha:\mathsf{Size}\), write
\({\downarrow \alpha} = (\beta:\mathsf{Size})\times (\beta<\alpha)\) for the type of sizes
strictly below \(\alpha\). We assume that the size structure has the
upper-bound and successor-like operations specified in Section 5.2.

A size-indexed approximation consists of a family
\(D : \mathsf{Size} \to \mathcal{U}\) together with operations
\(\tau_{\beta,\alpha} : T_\Sigma (D_\beta) \to D_\alpha\) whenever \(\beta<\alpha\). The map
\(\tau_{\beta,\alpha}\) interprets a term whose variables lie at stage \(\beta\) as
 an element constructed at the later stage \(\alpha\).

The stage \(D_\alpha\) is defined by well-founded recursion from the stages
strictly below \(\alpha\). Before quotienting, its elements are pairs
\((\beta,t) : (\beta:\mathsf{Size})\times (\beta<\alpha)\times T_\Sigma (D_\beta).\)

Such a pair records a lower stage \(\beta<\alpha\) together with a term \(t\)
whose variables have already been constructed at that stage. The next
approximation is obtained by quotienting these pairs by a relation
\(R_\alpha\):
\(D_\alpha = \left( (\beta:\mathsf{Size})\times (\beta<\alpha)\times T_\Sigma (D_\beta) \right)\big/ R_\alpha.\)

We write the equivalence class of \((\beta,t)\) as \([\beta,t]_\alpha : D_\alpha.\)

The relation \(R_\alpha\) is generated by three kinds of identification. The
first imposes the specified equations. For every \(e:E\), every \(\beta<\alpha\),
and every substitution \(\rho:V (e)\to D_\beta\), it identifies
\([\beta,T_\Sigma (\rho) (l (e))]_\alpha = [\beta,T_\Sigma (\rho) (r (e))]_\alpha.\)

The second identification expresses the unit law for the free monad. A
term inserted as a variable after already being interpreted at a lower
stage is identified with the original term. Schematically, for
\(\gamma<\beta<\alpha\) and \(t:T_\Sigma (D_\gamma)\), it identifies
\([\beta,\eta (\tau_{\gamma,\beta} (t))]_\alpha = [\gamma,t]_\alpha.\)

The third identification expresses compatibility with the multiplication
of the free monad. Constructing a term of terms and then interpreting
its inner terms is identified with the corresponding flattened term.
Equivalently, applying an operation before or after transporting its
immediate subterms to a later stage gives the same element of the
quotient.

These three classes of relations correspond respectively to the
equations of the presentation, the unit law, and the multiplication law.
Their inclusion ensures that the colimit will satisfy the specified
equations and carry an algebra for the free monad \(T_\Sigma\), and
hence an algebra for \(S_\Sigma\).

For every \(\beta<\alpha\), the structure operation is defined by
\(\tau_{\beta,\alpha} (t) = [\beta,t]_\alpha.\)

If \(\alpha<\beta\), the transition map \(\delta_{\alpha,\beta} : D_\alpha \to D_\beta\) is
obtained by inserting an element as a variable:
\(\delta_{\alpha,\beta} = \tau_{\alpha,\beta}\circ\eta.\)

On representatives, this map simply regards a term constructed below
\(\alpha\) as the same term occurring below the later stage \(\beta\):
\(\delta_{\alpha,\beta} ([\gamma,t]_\alpha) = [\gamma,t]_\beta.\)

The relations \(R_\alpha\) ensure that these transition maps are well-defined
and coherent. In particular, whenever \(\alpha<\beta<\gamma\), we have
\(\delta_{\beta,\gamma}\circ\delta_{\alpha,\beta} = \delta_{\alpha,\gamma}.\)

Thus \((D,\delta)\) forms a diagram indexed by \(\mathsf{Size}\). The
use of well-founded recursion separates the construction of a stage from
the quotients being formed at that stage and avoids treating quotient
formation as a strictly positive operation.

This construction is \emph{inflationary} because a stage contains terms
formed over all strictly smaller stages rather than only terms formed
over one immediately preceding stage. At a limit-like size, the
construction can therefore use everything constructed below that size
without requiring a separate successor or limit clause.

\subsection{4.2 The colimit
construction}\label{the-colimit-construction}

The candidate QW-type is the colimit of the diagram of approximants:
\(Q = \mathrm{colim}_{\alpha:\mathsf{Size}} D_\alpha.\)

Write \(\nu_\alpha : D_\alpha \to Q\) for the canonical colimit maps. These
satisfy \(\nu_\beta\circ\delta_{\alpha,\beta} = \nu_\alpha\) whenever \(\alpha<\beta\).

Concretely, the colimit can be presented as a quotient of the disjoint
union of the stages:
\(Q = \left( (\alpha:\mathsf{Size})\times D_\alpha \right)\big/{\sim},\) where
\((\alpha,x)\sim (\beta,y)\) when \(x\) and \(y\) become equal after being
transported to a sufficiently large common stage. The directedness
properties of \(\mathsf{Size}\) ensure that such common extensions can
be found.

The relations imposed at each stage pass to the colimit. In particular,
whenever a substitution \(\rho:V (e)\to Q\) can be lifted to a common
stage, the two interpretations of \(l (e)\) and \(r (e)\) are identified
at a later stage and hence become equal in \(Q\). The unit and
multiplication relations similarly ensure that the interpretation of
terms is independent of the stages and representatives used to construct
them.

To make \(Q\) into an \(S_\Sigma\)-algebra, we require a map
\(\mathsf{intro} : S_\Sigma (Q) \to Q.\)

At each stage \(D_\alpha\), one can add a new constructor layer at a
sufficiently later stage. These stagewise operations induce a map
\(\mu_D : \mathrm{colim}_\alpha S_\Sigma (D_\alpha) \to Q.\)

The map \(\mu_D\) is not yet the desired constructor, because its domain
consists of operations whose entire family of arguments already lies in
one stage. An arbitrary element of \(S_\Sigma (Q)\) consists instead of
an operation whose arguments are given only as elements of the colimit.
The remaining problem is therefore to lift such a family to a common
stage.

Once the algebra structure and equation proofs have been constructed,
the recursor into any other QW-algebra is defined stage by stage by
well-founded recursion. The compatibility relations make the resulting
family of maps \(D_\alpha\to X\) into a cocone. The universal property of the
colimit then produces a map \(Q\to X\), and a second well-founded
induction establishes its uniqueness. Thus the well-founded size
structure controls both the formation of the approximants and the
termination of the recursion defining the universal morphism.

\subsection{4.3 The cocontinuity
bottleneck}\label{the-cocontinuity-bottleneck}

For any diagram \(D\), functoriality gives a canonical map

\[
\kappa_D :
\mathrm{colim}_\alpha S_\Sigma (D_\alpha)
\longrightarrow
S_\Sigma \left(\mathrm{colim}_\alpha D_\alpha\right).
\]

It is induced by the cocone
\(S_\Sigma (\nu_\alpha) : S_\Sigma (D_\alpha) \to S_\Sigma (Q).\)

An element of the domain of \(\kappa_D\) consists of an operation
\(a:A\) and a family of arguments \(B (a)\to D_\alpha\) lying in a common
stage \(\alpha\). An element of its codomain consists of an operation \(a:A\)
and a family \(b : B (a) \to Q\) whose members may be represented at
different stages. Surjectivity of \(\kappa_D\) therefore states that
every such family can be lifted to some common stage.

If \(\kappa_D\) is an isomorphism, the constructor on the colimit is
obtained as the composite

\[
S_\Sigma (Q)
\xrightarrow{\kappa_D^{-1}}
\mathrm{colim}_\alpha S_\Sigma (D_\alpha)
\xrightarrow{\mu_D}
Q.
\]

For the construction of the algebra structure, a full inverse is
stronger than necessary. It is enough to have a section

\[
\sigma_D :
S_\Sigma (Q)
\longrightarrow
\mathrm{colim}_\alpha S_\Sigma (D_\alpha)
\] satisfying \(\kappa_D\circ\sigma_D = \mathsf{id}_{S_\Sigma (Q)}\).

The constructor can then be defined by
\(\mathsf{intro} = \mu_D\circ\sigma_D.\)

The equation system introduces an analogous lifting requirement. To
instantiate an equation \(e:E\) in \(Q\), a substitution
\(\rho : V (e) \to Q\) must be represented at a common stage. It is
convenient to package both kinds of lifting into a single polynomial
functor. Define \(C : A+E \to \mathcal{U}\) by
\(C (\mathsf{inl} (a)) = B (a)\) and \(C (\mathsf{inr} (e)) = V (e),\)
and let \(S_{\Sigma,E} (X) = (c:A+E)\times (C (c)\to X).\)

A section of the canonical map

\[
\kappa_D^{\Sigma,E} :
\mathrm{colim}_\alpha S_{\Sigma,E} (D_\alpha)
\longrightarrow
S_{\Sigma,E} \left(\mathrm{colim}_\alpha D_\alpha\right)
\] simultaneously supplies common-stage lifts for constructor arguments
and for substitutions into equations. This is the precise cocontinuity
property needed by the inflationary construction.

Fiore, Pitts and Steenkamp use WISC to construct a type of sizes
sufficiently rich for the relevant polynomial functor to preserve all
size-indexed colimits \citep{fiore2022-quotient-inductive}. For that
size structure, the canonical map \(\kappa_D^{\Sigma,E}\) is an
isomorphism for every size-indexed diagram \(D\). Its inverse provides
all the common-stage lifts required in the construction.

Our argument recovers only the part of this theorem used for the
particular inflationary diagram. We do not prove that the polynomial
functor preserves arbitrary size-indexed colimits. Instead, we construct
directly a section

\[
\sigma_D^{\Sigma,E} :
S_{\Sigma,E} (Q)
\longrightarrow
\mathrm{colim}_\alpha S_{\Sigma,E} (D_\alpha)
\] of \(\kappa_D^{\Sigma,E}\) for the diagram of approximants. The
construction of this section from quotient-invariant ranks is the
central technical task of the following sections.

\section{5. Extensional plump
ordinals}\label{extensional-plump-ordinals}

The inflationary construction requires a notion of stage with three
interacting properties. Stages must support well-founded recursion,
admit upper bounds for the arities occurring in the QW-presentation, and
assign equal stages to extensionally equivalent families. Ordinary
inductively generated size types provide the first two properties but
not, in general, the third. We therefore replace them with extensional
plump ordinals.

\subsection{5.1 Constructive ordinal
notions}\label{constructive-ordinal-notions}

The simplest type-theoretic representation of sizes is a W-type. Given a
signature \(\Sigma=(A,B)\) of supremum constructors, consider

\[
W_\Sigma = \mathsf{W} (a:A).B (a).
\]

An element of \(W_\Sigma\) is a well-founded tree \(\mathsf{sup} (a,f)\)
with \(a:A\) and \(f:B (a)\to W_\Sigma\). By choosing suitable elements
of \(A\), the W-type can include nullary, binary, and signature-specific
supremum operations.

The inductive structure of \(W_\Sigma\) supplies a well-founded relation
suitable for recursion. The plump ordering additionally ensures that
\(\mathsf{sup} (a,f)\) is a strict upper bound for the family \(f\).
Fiore, Pitts and Steenkamp use such plump W-types as their size
structures \citep{fiore2022-quotient-inductive}, following the
constructive theory of plump ordinals introduced by Taylor
\citep{taylor1993-intuinistic-sets-ordinals}. Directed variants of the
plump ordering provide finite joins and related closure properties
\citep{gratzer2024-universe-grothendieck, gratzer2022-directed-plump-order}.

The equality of a W-type is intensional. Two trees may describe the same
well-founded order while remaining unequal because they were built using
different constructors or different enumerations of their immediate
subtrees. In particular, a family \(f:B (a)\to W_\Sigma\) and a
reindexing \(f\circ\pi\) generally determine distinct W-trees even when
\(\pi\) is an isomorphism.

An extensional well-founded order instead identifies an ordinal by the
order type of its predecessors. Writing

\[
{\downarrow\alpha} = (\gamma:Z)\times (\gamma<\alpha),
\]

extensionality asserts that ordinals with the same strict lower segment
are equal. In one standard formulation,

\[
\left(
(\gamma:Z)\to
(\gamma<\alpha\leftrightarrow\gamma<\beta)
\right)
\longrightarrow
\alpha=\beta.
\]

Extensional well-founded orders provide representative-independent
ordinal values, but they do not automatically contain upper bounds for
all the families required by a given infinitary signature. Plumpness
supplies these upper bounds and ensures that they interact
constructively with the order. Our size structure combines this closure
with extensional equality.

\subsection{5.2 Extensional plump ordinal
bases}\label{extensional-plump-ordinal-bases}

Fix a supremum signature \(\Sigma=(A,B)\). It consists of a type
\(A:\mathcal{U}\) and a family \(B:A\to\mathcal{U}\). The arities \(B (a)\) are chosen to
include the constructor arities and equation-variable contexts relevant
to the QW-presentation. We additionally require enough nullary and
binary structure to construct a least element and finite upper bounds.

A plump algebra for \(\Sigma=(A,B)\) consists first of a type \(Z\) equipped
with two relations

\[
{<},{\leq}:Z\to Z\to\mathsf{Prop}.
\]

The strict relation \(<\) represents strict comparison of stages, while
\(\leq\) represents non-strict comparison. The algebra has a supremum
operation

\[
\mathsf{sup}: (a:A)\times (B (a)\to Z) \to Z.
\]

For \(a:A\) and \(f:B (a)\to Z\), we write its value as
\(\mathsf{sup} (a,f).\)

This operation satisfies two plumpness inequalities. First, if every
member of the family is strictly below \(\alpha\), then the generated
supremum is non-strictly below \(\alpha\):

\[
\left(
(\xi:B (a))\to f (\xi)<\alpha
\right)
\longrightarrow
\mathsf{sup} (a,f)\leq\alpha.
\]

Second, anything non-strictly below one member of the family is strictly
below the generated supremum:

\[
\alpha\leq f (\xi) \longrightarrow \alpha<\mathsf{sup} (a,f).
\]

Taking \(\alpha=f (\xi)\) in the second inequality shows that every member
is strictly below the generated supremum: \(f (\xi)<\mathsf{sup} (a,f).\)

Thus \(\mathsf{sup} (a,f)\) is a strict upper bound rather than merely a
least upper bound in the usual non-strict sense. The first inequality
expresses its minimality relative to other strict upper bounds.

The relations are required to satisfy the expected composition laws. The
relation \(\leq\) is reflexive and transitive: \(\alpha\leq\alpha\) and

\[
\alpha\leq\beta
\;\land\;
\beta\leq\gamma
\longrightarrow
\alpha\leq\gamma.
\]

Strict comparison composes with both strict and non-strict comparison:
\(\alpha<\beta \;\land\; \beta\leq\gamma \longrightarrow \alpha<\gamma\),
\(\alpha\leq\beta \;\land\; \beta<\gamma \longrightarrow \alpha<\gamma\),
and
\(\alpha<\beta \;\land\; \beta<\gamma \longrightarrow \alpha<\gamma\).

Strict comparison also implies non-strict comparison:
\(\alpha<\beta \longrightarrow \alpha\leq\beta.\)

The algebra contains a binary strict upper-bound operation

\[
{-}\mathbin{\vee^Z}{-}:Z\to Z\to Z.
\]

It satisfies

\[
\alpha<\alpha\vee^Z\beta
\qquad\text{and}\qquad
\beta<\alpha\vee^Z\beta.
\]

It is minimal among common strict upper bounds in the sense that

\[
\alpha<\gamma
\;\land\;
\beta<\gamma
\longrightarrow
\alpha\vee^Z\beta\leq\gamma.
\]

We also require the comparison
\(\beta\vee^Z\alpha \leq \alpha\vee^Z\beta.\)

Applying the same law with \(\alpha\) and \(\beta\) exchanged gives the
reverse comparison. Antisymmetry will therefore imply that the binary
bound is commutative.

Finally, the algebra contains a least element \(\bot^Z:Z\) satisfying
\(\bot^Z\leq\alpha\) for every \(\alpha:Z\). The binary operation and
least element allow finitely many auxiliary bounds to be combined with
the signature-specific suprema.

A plump algebra is \emph{extensional} when its non-strict comparison is
antisymmetric:

\[
\alpha\leq\beta
\;\land\;
\beta\leq\alpha
\longrightarrow
\alpha=\beta.
\]

In the intended construction, \(\alpha\leq\beta\) expresses inclusion,
or simulation, of the strict predecessor structure of \(\alpha\) by that
of \(\beta\). Equality of strict lower segments therefore gives
comparisons in both directions. Antisymmetry then yields the
extensionality principle

\[
\left(
(\gamma:Z)\to
(\gamma<\alpha\leftrightarrow\gamma<\beta)
\right)
\longrightarrow
\alpha=\beta.
\]

The Agda implementation records antisymmetry rather than repeating this
lower-segment formulation because the connection between \(\leq\) and
predecessor inclusion is established separately.

We write \(\mathsf{ExtPlumpAlg} (\Sigma)\) for the type of extensional
plump algebras over \(\Sigma=(A,B)\). We require an initial algebra for
extensional plump ordinals:

\[
Z^e_\Sigma : \mathsf{ExtPlumpAlg} (\Sigma).
\]

Initiality provides the associated recursion and induction principles
and makes the ordinal structure canonical for the chosen arities. We
refer to \(Z^e_\Sigma\) as the \emph{extensional plump ordinal basis}
generated by \(\Sigma\). Thus the main construction assumes an initial
algebra for extensional plump ordinals, not merely an extensional plump
algebra.

The assumption used in the main construction can be packaged
universe-polymorphically as

\[
\mathsf{ExtensionalPlumpOrdinals}
=
(A:\mathcal{U})\to
(B:A\to\mathcal{U})\to
\mathsf{ExtPlumpAlg} (A,B),
\] where each chosen algebra is understood to be initial. In the
formalisation, universe levels are made explicit so that the
construction can be instantiated for signatures at different levels.

For a fixed QW-presentation \((\Sigma,E,V,l,r)\), we do not require
closure under every small type. It is enough to choose \(\Sigma=(A,B)\) whose
arities include the types \(B (a)\) for \(a:A\) and \(V (e)\) for
\(e:E\), together with nullary and binary arities. This is the sense in
which the main theorem is relative to a sufficiently closed extensional
plump ordinal basis.

\subsection{5.3 Extensional supremum
comparison}\label{extensional-supremum-comparison}

The following lemma packages the plump-order calculation used throughout
the rank argument.

\begin{quote}
\textbf{Lemma 5.1 (Extensional supremum comparison).} Let \(f:I\to Z\)
and \(g:J\to Z\) be families supported by supremum constructors
\(a,b:A\). If every \(f (\xi)\) equals some \(g (\eta)\), then
\(\mathsf{sup} (a,f)\leq\mathsf{sup} (b,g)\). If the two families
have the same image, their suprema are equal.
\end{quote}

For each \(\xi:I\), choose the given \(\eta:J\) with \(f (\xi)=g (\eta)\).
Reflexivity and the second plumpness inequality give
\(f (\xi)=g (\eta)<\mathsf{sup} (b,g)\). The first plumpness inequality
then yields the stated comparison. Applying the argument in both
directions and using antisymmetry proves equality. The same proof works
when image containment is propositionally truncated, because the target
inequalities are propositions.

In particular, reindexing by an isomorphism does not change a supremum:
\(\mathsf{sup} (a,f) = \mathsf{sup} (a,f\circ\pi)\) for
\(\pi:I\cong I\). This conclusion generally fails for intensional
W-type sizes: the two expressions can be different trees even when they
present the same predecessor ranks. Extensionality is therefore what
allows a plump rank to descend through permutation and image equations.

\section{6. Depth-preserving
signatures}\label{depth-preserving-signatures}

We now identify the class of QW-presentations for which the ordinal rank
of a raw term is invariant under the specified equations. The formal
definition is semantic: an equation must preserve rank under every
substitution of raw terms for its variables. We subsequently give a
stronger syntactic condition that makes this property easy to verify in
examples.

Fix a signature \(\Sigma=(A,B)\) and an equation system \((E,\Xi)\). For
each \(e:E\), the equation \(\Xi (e)\) consists of a variable type
\(V_e\) and two terms \(l_e,r_e:T_\Sigma (V_e).\)

Let \(T\) denote the W-type of closed raw terms over \(\Sigma\). Let
\(Z\) be an extensional plump ordinal basis for the same arities.

\subsection{6.1 Rank and depth preservation}\label{rank-and-depth-preservation}

The W-type \(T\) is generated by a constructor
\(\mathsf{sup}^T: (a:A)\times (B (a)\to T) \to T.\)

The corresponding extensional plump ordinal has a supremum operation
\(\mathsf{sup}^Z: (a:A)\times (B (a)\to Z) \to Z.\)

Initiality of the W-type gives a rank function \(\iota^Z:T\to Z\)
defined recursively by

\[
\iota^Z \left(\mathsf{sup}^T (a,f)\right)
=
\mathsf{sup}^Z \left(a,\lambda \xi.\iota^Z (f (\xi))\right).
\]

Thus the rank of a raw term is obtained by replacing every term
constructor with the corresponding ordinal supremum constructor. The
plumpness inequalities ensure that every immediate subterm has rank
strictly below the rank of the whole term.

Open terms in the free algebra are generated by variables and constructor
applications: \(\mathsf{var}^E:V\to T_\Sigma (V)\) and

\[
\mathsf{sup}^E:
(a:A)\times (B (a)\to T_\Sigma (V))
\to
T_\Sigma (V).
\]


Given a substitution \(\rho:V\to T,\) the interpretation of an
expression as a raw term is defined recursively by
\(\mathsf{assign}_T (\rho,\mathsf{var}^E (v)) = \rho (v)\) and

\[
\mathsf{assign}_T \left(\rho,\mathsf{sup}^E (a,f)\right)
=
\mathsf{sup}^T \left(a,\lambda \xi.\mathsf{assign}_T (\rho,f (\xi))\right).
\]

Equivalently, one may define the rank of an open term relative to
a rank valuation \(r:V\to Z.\)

Write this rank as \(\mathsf{rank}_r:T_\Sigma (V)\to Z.\)

It is given by \(\mathsf{rank}_r (\mathsf{var}^E (v)) = r (v)\) and

\[
\mathsf{rank}_r \left(\mathsf{sup}^E (a,f)\right)
=
\mathsf{sup}^Z \left(a,\lambda \xi.\mathsf{rank}_r (f (\xi))\right).
\]

The rank of a substituted expression is then

\[
\iota^Z (\mathsf{assign}_T (\rho,e))
=
\mathsf{rank}_{\iota^Z\circ\rho} (e).
\]

Assigning \(\bot^Z\) to every variable gives the corresponding bottom-rank
special case.

Because \(Z\) is extensional, depth preservation can be stated directly
as equality of ranks.

An equation \(e:E\) is \emph{depth-preserving} when, for every
substitution \(\rho:V_e\to T\), its two sides have equal ranks:

\[
\iota^Z (\mathsf{assign}_T (\rho,l_e))
=
\iota^Z (\mathsf{assign}_T (\rho,r_e)).
\]

A QW-signature is depth-preserving when every one of its equations is
depth-preserving:

\[
\mathsf{DepthPreserving} (\Sigma,E,\Xi)
=
(e:E)\to
(\rho:V_e\to T)\to
\left(
\iota^Z (\mathsf{assign}_T (\rho,l_e))
=
\iota^Z (\mathsf{assign}_T (\rho,r_e))
\right).
\]

This is the formal definition used in the construction. It is semantic
because it quantifies over substitutions by arbitrary raw terms and
compares their interpretations in the rank algebra \(Z\). It is
nevertheless a property of the syntactic presentation and can often be
established by a structural analysis of the two sides.

For signatures whose ordinal supremum depends only on the collection of
predecessor ranks, a useful proof principle is available. It is
enough that the two sides have the same overall depth and that, for each
variable, the collection of depths at which it occurs is the same on
both sides. Reordering or extensionally reindexing occurrences is
permitted. The number and order of occurrences need not be
definitionally equal when the extensional supremum identifies the
resulting predecessor collections.

This proof principle is stronger than the semantic definition. There
may be equations whose ranks agree for algebraic reasons even though
their variable occurrences do not match in this direct way. The semantic
formulation is therefore the appropriate general definition, while
occurrence-depth matching is a convenient proof principle for examples.

\subsection{6.2 Descent through the
quotient}\label{descent-through-the-quotient}

Let \(\approx\) be the congruence on \(T\) generated by all substitution
instances of the equations. Thus, for every \(e:E\) and
\(\rho:V_e\to T\), it contains
\(\mathsf{assign}_T (\rho,l_e) \approx \mathsf{assign}_T (\rho,r_e).\)

It is closed under reflexivity, symmetry, transitivity, and applications
of the constructors of \(\Sigma\).

Suppose that the signature is depth-preserving. For every generating
equation, depth preservation gives

\[
\iota^Z (\mathsf{assign}_T (\rho,l_e))
=
\iota^Z (\mathsf{assign}_T (\rho,r_e)).
\]

The rank function also respects the equivalence closure of the relation.
The reflexive, symmetric, and transitive cases follow immediately from
the corresponding properties of equality.

For the congruence case, suppose that \(f,g:B (a)\to T\) satisfy
\(f (\xi)\approx g (\xi)\) for every \(\xi:B (a)\). By induction on the
derivations of these relations, we have
\(\iota^Z (f (\xi)) = \iota^Z (g (\xi))\) for every \(\xi:B (a)\). Function
extensionality gives
\((\lambda \xi.\iota^Z (f (\xi))) = (\lambda \xi.\iota^Z (g (\xi))),\) and hence

\[
\mathsf{sup}^Z (a,\lambda \xi.\iota^Z (f (\xi)))
=
\mathsf{sup}^Z (a,\lambda \xi.\iota^Z (g (\xi))).
\]

By the recursive definition of \(\iota^Z\), this is precisely
\(\iota^Z (\mathsf{sup}^T (a,f)) = \iota^Z (\mathsf{sup}^T (a,g)).\)

It follows by induction on the generated congruence that
\(t\approx u \longrightarrow \iota^Z (t)=\iota^Z (u).\)

Let \(q:T\to T/{\approx}\) be the quotient map. The quotient elimination
principle therefore defines a rank function
\(\overline{\iota}^Z:T/{\approx}\to Z\) satisfying
\(\overline{\iota}^Z (q (t)) = \iota^Z (t).\)

This rank is independent of the representative chosen for an equivalence
class. It is the canonical stage assignment that is unavailable when one
quotients raw terms without a depth-preservation argument.

The inflationary construction uses quotients at every stage rather than
only the single quotient \(T/{\approx}\). The same proof applies
stagewise. The equation relations preserve rank by depth preservation,
while the unit and multiplication relations preserve rank by the algebra
laws for the recursive interpretation. Consequently, each approximant
\(D_\alpha\) inherits a rank function compatible with the transition maps.
These compatible functions induce a rank \(\mathsf{rank}_Q:Q\to Z\) on
the colimit of the approximants.

\subsection{6.3 Examples}\label{examples}

For mobiles, a substitution \(\rho:\mathbb{N}\to T\) turns a permutation
equation into suprema of \(\iota^Z\circ\rho\) and
\(\iota^Z\circ\rho\circ\pi\). These families have the same image, so
Lemma 5.1 makes their ranks equal. The same argument proves depth
preservation for the image-based presentation.

For the cumulative hierarchy, an equation is induced by \(f:A\to X\) and
\(g:B\to X\) with the same image. After any substitution
\(\rho:X\to T\), the families \(\iota^Z\circ\rho\circ f\) and
\(\iota^Z\circ\rho\circ g\) again have the same image. Lemma 5.1
therefore identifies the two ranks, even though their indexing types
differ.

For finite multisets, the adjacent-transposition equation places its
sole recursive variable beneath two unary constructors on both sides.
The occurrence-depth criterion of Section 6.1 applies directly, so the
presentation is depth-preserving.

Associativity and unit are non-examples:
\((x\star y)\star z=x\star (y\star z)\) changes the depths of \(x\) and
\(z\), while \(e\star x=x\) changes the depth of \(x\). Their finitary
QW-types remain constructible by Section 8.2; depth preservation is
sufficient for the rank construction, not necessary for existence.

\section{7. Cocontinuity from canonical
ranks}\label{cocontinuity-from-canonical-ranks}

We now prove the central technical result. For the diagram of
rank-bounded approximants, exact-rank lifting constructs a section of
the canonical map from the colimit of the polynomial functor to the
polynomial functor of the colimit.

The formal proof is carried out in extensional dependent type theory
with path elimination, function extensionality, propositional
extensionality, unique choice, and set quotients. We assume a signature
\(\Sigma=(A,B)\), an equation system \((E,\Xi)\), a proof that the
signature is depth-preserving, and an extensional plump ordinal basis
\(Z\) for \(\Sigma\). No instance of WISC or of general choice is used in
this argument.

\subsection{7.1 The canonical colimit
map}\label{the-canonical-colimit-map}

Let \(T\) be the W-type of raw terms over \(\Sigma\), and let
\(\iota^Z:T\to Z\) be its ordinal-valued rank. For each \(\alpha:Z\),
define the type of raw terms bounded by \(\alpha\) by
\(S_0 (\alpha) = (t:T)\times\bigl(\iota^Z (t)\leq\alpha\bigr).\)

The equation system generates a setoid \(\widetilde{S} (\alpha)\) on
\(S_0 (\alpha)\). Its equivalence relation contains the specified
equations and is closed under constructor congruence, reflexivity,
symmetry, and transitivity. Write \(\widetilde{S}/ (\alpha)\) for the
set quotient of this setoid.

Whenever \(\alpha\leq\beta\), a bounded term at stage \(\alpha\) is also
bounded at stage \(\beta\). This gives a weakening map

\[
\mathsf{weaken}_{\alpha,\beta}:
\widetilde{S}/ (\alpha)
\to
\widetilde{S}/ (\beta).
\]

These maps are functorial, so the assignment
\(\alpha\longmapsto\widetilde{S}/ (\alpha)\) defines the diagram
\(D:(Z,\leq)\to\mathsf{Set}.\)

Its colimit is
\(\mathrm{colim}\,D = \mathrm{colim}_{\alpha:Z} \widetilde{S}/ (\alpha).\)

A representative of an element of this colimit is a pair
\((\alpha,\widehat{t})\) with \(\widehat{t}:\widetilde{S}/ (\alpha)\).
We write its equivalence class as \([\alpha,\widehat{t}]_D:\mathrm{colim}\,D.\)

The colimit relation identifies elements equal at the same stage and
elements related by weakening to a larger stage. It is then closed under
symmetry and transitivity.

Let \(S\) be the polynomial functor associated with the container
\(\Sigma=(A,B)\): \(S (X) = (a:A)\times (B (a)\to X).\)

Applying \(S\) pointwise to the diagram gives the diagram
\(S\circ D:(Z,\leq)\to\mathsf{Set}.\)

Its colimit is
\(\mathrm{colim}(S\circ D) = \mathrm{colim}_{\alpha:Z} S (\widetilde{S}/ (\alpha)).\)

An element of \(\mathrm{colim}(S\circ D)\) is represented by a stage
\(\alpha:Z\), a constructor \(a:A\), and a family
\(\widehat{f}:B (a)\to\widetilde{S}/ (\alpha).\)

We write this representative as \([\alpha,a,\widehat{f}]_{SD}.\)

The colimit cocone for \(D\) induces the canonical map
\(\varphi: \mathrm{colim}(S\circ D) \longrightarrow S (\mathrm{colim}\,D).\)

On representatives, it is defined by

\[
\varphi ([\alpha,a,\widehat{f}]_{SD})
=
\left(
a,\lambda \xi.[\alpha,\widehat{f} (\xi)]_D
\right).
\]

Thus \(\varphi\) forgets the common stage \(\alpha\) and maps each
argument into the colimit. Functoriality of weakening ensures that this
definition respects the colimit relation.

For the stage construction used here, the equation relation has already
been imposed inside each quotient \(\widetilde{S}/ (\alpha)\).
Consequently, the cocontinuity problem concerns the underlying
polynomial functor \(S\) alone. No additional polynomial containing the
equation-variable contexts is required in the formal theorem.

The required result is a section
\(\sigma_D: S (\mathrm{colim}\,D) \longrightarrow \mathrm{colim}(S\circ D).\)

We must prove \(\varphi\circ\sigma_D = \mathsf{id}_{S (\mathrm{colim}\,D)}\), which
is precisely the common-stage lifting needed to define constructors on
the colimit.

\subsection{7.2 Exact-rank lifting}\label{exact-rank-lifting}

We first construct a canonical rank on every stage quotient. Define
\(\mathsf{rank}_0: S_0 (\alpha)\to Z\) by
\(\mathsf{rank}_0 (t,p) = \iota^Z (t).\)

Depth preservation implies that equivalent bounded terms have equal
ranks. The constructor-congruence case follows recursively from
congruence of \(\mathsf{sup}^Z\). For a generating equation, depth
preservation gives equality of the two ranks.

The quotient recursion principle therefore yields
\(\mathsf{rank}_\alpha: \widetilde{S}/ (\alpha)\to Z.\)

For a represented bounded term, it satisfies
\(\mathsf{rank}_\alpha ([(t,p)]) = \iota^Z (t).\)

The bound carried by the representative gives
\(\mathsf{rank}_\alpha (\widehat{t}) \leq \alpha\) for every
\(\widehat{t}:\widetilde{S}/ (\alpha)\). Moreover, weakening does not
change rank:
\(\mathsf{rank}_\beta (\mathsf{weaken}_{\alpha,\beta} (\widehat{t})) = \mathsf{rank}_\alpha (\widehat{t})\)
whenever \(\alpha\leq\beta\).

These stagewise ranks respect the colimit relation. They therefore
induce a canonical colimit rank \(\mathsf{rank}_C:\mathrm{colim}\,D\to Z\)
satisfying

\[
\mathsf{rank}_C ([\alpha,\widehat{t}]_D)
=
\mathsf{rank}_\alpha (\widehat{t}).
\]

The crucial strengthening is that every colimit element can be
represented at exactly its canonical rank. The rank descent and the
dependent transport needed for this statement are collected in one
lemma.

\begin{quote}
\textbf{Lemma 7.1 (Canonical exact-rank lifting).} There are functions
\[
\mathsf{rank}_C:\mathrm{colim}\,D\to Z
\] and \[
\mathsf{lift}_C:
(x:\mathrm{colim}\,D)\to\widetilde{S}/ (\mathsf{rank}_C (x))
\] such that \([\mathsf{rank}_C (x),\mathsf{lift}_C (x)]_D=x\). They are
invariant under the stage equations and compatible with weakening.
\end{quote}

To prove the lemma, first define the lift on a bounded raw term. For a
bounded raw term \((t,p):S_0 (\alpha),\) define its exact-rank
representative by

\[
\mathsf{lift}_0 (t,p)
=
\left(
t,
\mathsf{refl}_{\leq} (\iota^Z (t))
\right)
:
S_0 (\iota^Z (t)).
\]

Depth preservation makes equivalent raw terms have equal ranks. After
transport along that equality, proof irrelevance for order witnesses
makes their exact-rank representatives equal. Quotient recursion
therefore gives

\[
\mathsf{lift}_{\approx} (\widehat{t})
:
\widetilde{S}/ (\mathsf{rank}_\alpha (\widehat{t})).
\]

Compatibility with weakening lets this operation descend through the
colimit relation, proving Lemma 7.1. No representative is chosen: the
target stage is determined intrinsically by the quotient-invariant rank.

We now define the section of the canonical colimit map. Let
\((a,f):S (\mathrm{colim}\,D),\) so that \(f:B (a)\to\mathrm{colim}\,D.\)

Lift each argument canonically: \(\mu (\xi)=\mathsf{rank}_C (f (\xi))\) and
\(\widehat{g} (\xi)=\mathsf{lift}_C (f (\xi))\), where
\(\widehat{g} (\xi):\widetilde{S}/ (\mu (\xi))\).

Use the plump supremum to define a common stage
\(\alpha = \mathsf{sup}^Z (a,\mu).\)

For every \(\xi:B (a)\), the plumpness inequalities give
\(\mu (\xi)<\alpha\) and hence \(\mu (\xi)\leq\alpha.\)

We may therefore weaken each exact-rank representative to stage
\(\alpha\):

\[
\widehat{x} (\xi)
=
\mathsf{weaken}_{\mu (\xi),\alpha} (\widehat{g} (\xi))
:
\widetilde{S}/ (\alpha).
\]

The family \(\widehat{x}\) lies entirely at the common stage \(\alpha\).
We define \(\sigma_D (a,f) = [\alpha,a,\widehat{x}]_{SD}.\)

This construction contains no arbitrary representative choices. The
ranks, exact-rank representatives, and common supremum stage are all
canonical.

\subsection{7.3 The section theorem}\label{the-section-theorem}

Let \((a,f):S (\mathrm{colim}\,D)\). By definition,

\[
\varphi (\sigma_D (a,f))
=
\left(
 a,
 \lambda \xi.[\alpha,\widehat{x} (\xi)]_D
\right).
\]

For each \(\xi:B (a)\), the element \(\widehat{x} (\xi)\) is the weakening
to \(\alpha\) of the exact-rank representative \(\widehat{g} (\xi)\). The
colimit relation therefore gives
\([\alpha,\widehat{x} (\xi)]_D = [\mu (\xi),\widehat{g} (\xi)]_D.\)

The exact-rank lifting theorem gives
\([\mu (\xi),\widehat{g} (\xi)]_D = f (\xi).\)

Combining these equalities pointwise and applying function
extensionality yields \(\varphi (\sigma_D (a,f)) = (a,f).\)

Hence \(\varphi\circ\sigma_D = \mathsf{id}_{S (\mathrm{colim}\,D)}.\)

This proves the required theorem.

\begin{quote}
\textbf{Theorem 7.2 (Cocontinuity from depth preservation).} Let
\(\Sigma=(A,B)\) be a QW-signature with a depth-preserving equation
system. Assume an extensional plump ordinal basis \(Z\) supporting the
arities \(B (a)\). For the diagram \(D\) of rank-bounded
quotient approximants, the canonical map \[
\varphi:
\mathrm{colim}(S\circ D)
\longrightarrow
 S (\mathrm{colim}\,D)
\] admits the section \(\sigma_D\) constructed above.
\end{quote}

The three hypotheses have distinct roles in the proof. Depth
preservation ensures that the raw rank is invariant under the generating
equations. Extensionality ensures that these rank equalities descend
through the stage quotients. Plumpness places all arguments of a constructor below the
common supremum stage.

\subsection{7.4 Cocontinuity and
existence}\label{cocontinuity-and-existence}

Theorem 7.2 is weaker than the WISC-based preservation theorem of Fiore,
Pitts and Steenkamp: it constructs a section only for the diagram
arising in the QW-construction, rather than an isomorphism for every
diagram indexed by the chosen sizes
\citep{fiore2022-quotient-inductive}. That restricted result is
sufficient.

Let \(Q = \mathrm{colim}\,D = \mathrm{colim}_{\alpha:Z} \widetilde{S}/ (\alpha).\)

The stage construction provides a map
\(\mu: \mathrm{colim}(S\circ D) \to Q\)
which adds one constructor layer. Using the section \(\sigma_D\) of the
canonical map, define
\(\mathsf{intro} = \mu\circ\sigma_D : S (Q)\to Q.\)

Because the specified equations have already been imposed in the stage
quotients, this algebra satisfies the equation system. The compatibility
of the stage relations ensures that the resulting interpretation is
independent of all stage representatives.

Given any other algebra \((X,a)\) satisfying the equations, define maps
\(r_\alpha: \widetilde{S}/ (\alpha)\to X\) by well-founded recursion on
\(\alpha\). The equation-satisfaction hypothesis for \(X\) makes each
\(r_\alpha\) well-defined on the stage quotient. The recursion equations
imply compatibility with weakening:
\(r_\beta\circ\mathsf{weaken}_{\alpha,\beta} = r_\alpha\) whenever
\(\alpha\leq\beta\). The colimit universal property therefore induces a
homomorphism \(r:Q\to X.\)

Well-founded induction on ranks shows that any algebra homomorphism
\(Q\to X\) agrees with the stagewise maps \(r_\alpha\). It follows that
\(r\) is unique. Hence \((Q,\mathsf{intro})\) is initial among algebras
satisfying the equation system.

Combining the cocontinuity theorem with the inflationary construction
yields the main result.

\begin{quote}
\textbf{Theorem 7.3 (Existence of depth-preserving QW-types).} Work in
extensional dependent type theory with W-types, set quotients, path
elimination, function extensionality, propositional extensionality, and
unique choice. Let \(\Sigma=(A,B)\) be a polynomial signature, and let
\((E,\Xi)\) be a depth-preserving system of equations over \(\Sigma\).
Suppose that there exists an initial extensional plump ordinal algebra
for \(\Sigma\). Then the QW-type presented by \((\Sigma,E,\Xi)\) can be
constructed from W-types and set quotients.
\end{quote}

Equivalently, every depth-preserving QW-signature whose constructor
arities are supported by an extensional plump ordinal basis has an
initial QW-algebra. The construction does not assume WISC or any general
choice principle beyond the existence of the specified ordinal basis.

\section{8. Beyond depth preservation}\label{beyond-depth-preservation}

The proof of Section 7 suggests a criterion more general than depth
preservation.

\begin{quote}
Inflationary iteration succeeds whenever the canonical map associated
with the relevant colimit admits a coherent section.
\end{quote}

For a diagram \(D:I\to\mathsf{Set}\) and a polynomial functor \(S\),
this canonical map is

\[
\kappa_D:
\mathrm{colim}_\alpha S (D_\alpha)
\longrightarrow
S \left(\mathrm{colim}_\alpha D_\alpha\right).
\]

A section

\[
\sigma_D:
S \left(\mathrm{colim}_\alpha D_\alpha\right)
\longrightarrow
\mathrm{colim}_\alpha S (D_\alpha)
\] satisfying \(\kappa_D\circ\sigma_D = \mathsf{id}\) assigns a
common-stage representative to every constructor application over the
colimit. Composing \(\sigma_D\) with the stagewise constructor gives an
algebra structure on the colimit.

The section must also be coherent with the transition maps, the stage
quotients, and the equations of the presentation. These coherence
properties are needed to define the universal morphism and prove that it
is an algebra homomorphism. The exact-rank construction of Section 7
verifies this coherence directly; a two-sided inverse to \(\kappa_D\) is
unnecessary.

Depth preservation is one mechanism for constructing this section. It
produces a canonical quotient-invariant rank, an exact-rank
representative, and a common stage obtained by applying the plump
supremum. Other mechanisms can provide the same data without preserving
the original structural depth.

\subsection{8.1 Normalising signatures}\label{normalising-signatures}

Suppose that raw terms admit a normalisation operation
\(\mathsf{nf}:T\to T.\)

For normalisation to yield a rank on the quotient, it must at least be
sound with respect to the equations: \(t\approx\mathsf{nf} (t).\)

It must also be invariant under the generated equivalence relation:
\(t\approx u \longrightarrow \mathsf{nf} (t)=\mathsf{nf} (u).\)

A stronger and more familiar condition is completeness:

\[
t\approx u
\quad\Longleftrightarrow\quad
\mathsf{nf} (t)=\mathsf{nf} (u).
\]

Idempotence is also desirable:
\(\mathsf{nf} (\mathsf{nf} (t)) = \mathsf{nf} (t).\)

Under these assumptions, define the rank of an equivalence class by
\(\mathsf{rank} ([t]) = \iota^Z (\mathsf{nf} (t)).\)

Invariance of normal forms makes this definition independent of the
representative \(t\). The term \(\mathsf{nf} (t)\) then provides a
preferred representative of the equivalence class at the stage
determined by its rank.

For use in the inflationary diagram, normalisation must additionally
commute with weakening. If \(t\) is represented at stage \(\alpha\) and
weakened to stage \(\beta\), normalising before or after weakening must
produce equivalent stage elements. This compatibility allows the
preferred representatives to descend through the colimit relation.

Given these properties, the construction of Section 7 can be repeated
with normal-form rank in place of structural rank. Each colimit element
is represented by its normal form at the rank of that normal form. The
ranks of the arguments of a constructor determine a common stage, and
the resulting family defines a section of \(\kappa_D\).

Normalisation can therefore handle equations that change syntactic
depth. Associativity and unit equations are typical examples for which
raw structural rank is not invariant but canonical normal forms may
still exist. The essential condition is not that the original equations
preserve rank, but that every equivalence class has a coherently chosen
representative whose rank is stable.

This argument does not require equations to be characterised exactly by
normalisation if a weaker result suffices. It is enough that equivalent
terms normalise to terms with equal ranks and that the normalised terms
represent the original equivalence classes. Equality of normal forms is
a convenient strong hypothesis because it supplies both properties
simultaneously.

\subsection{8.2 Finitary signatures}\label{finitary-signatures}

Finitary signatures admit a more direct argument. If every constructor
has only finitely many arguments, then a family of colimit elements
occurring in one constructor application can be lifted to a common stage
using finite directedness. No extensional plump ordinal rank is needed
to make this lift.

Constructively, the relevant notion is \emph{mere finiteness}. For a set
\(A\), define

\[
\mathsf{isFinite}^{\mathsf{p}} (A)
=
\exists n:\mathbb{N}.\,
\left\lVert
\mathsf{Fin} (n)\cong A
\right\rVert.
\]

The existential and the isomorphism are propositionally truncated. Thus
mere finiteness does not include a chosen cardinality or a chosen
enumeration of \(A\).

We also consider the type

\[
\mathsf{isFinite}' (A)
=
(n:\mathbb{N})\times
\left\lVert
\mathsf{Fin} (n)\cong A
\right\rVert.
\]

Although this type contains a natural number, that number is unique
whenever it exists. If \(\mathsf{Fin} (m)\cong A\) and
\(\mathsf{Fin} (n)\cong A\), the induced injections give \(m\leq n\) and
\(n\leq m\), so antisymmetry of the natural-number order gives \(m=n\).
This finite argument does not invoke the unrestricted
Cantor--Schröder--Bernstein theorem, which implies excluded middle
constructively \citep{pradic2019-cantor-bernstein}.

Consequently,
\((n:\mathbb{N})\times \left\lVert \mathsf{Fin} (n)\cong A \right\rVert\)
is a proposition. Unique choice therefore gives
\(\mathsf{isFinite}^{\mathsf{p}} (A) \longrightarrow \mathsf{isFinite}' (A)\).

This extracts the unique cardinality of \(A\), but it does not extract a
particular isomorphism \(\mathsf{Fin} (n)\cong A\) from its truncation.
For most of the cocontinuity argument, no such choice of enumeration is
needed. The desired common-stage property is itself a proposition, so
the truncated isomorphism can be eliminated while proving it.

Let \(D:I\to\mathsf{Set}\) be a filtered diagram and suppose that \(A\)
is merely finite. Given a function \(f:A\to\mathrm{colim}_\alpha D_\alpha,\) mere
finiteness reduces the problem propositionally to a family indexed by
\(\mathsf{Fin} (n)\). Finite directedness then gives a stage \(\alpha\) and a
function \(\widehat{f}:A\to D_\alpha\) such that
\(\nu_\alpha\circ\widehat{f} = f.\)

It follows that the canonical map
\(\mathrm{colim}_\alpha (A\to D_\alpha) \longrightarrow (A\to\mathrm{colim}_\alpha D_\alpha)\)
is surjective. Filteredness similarly proves injectivity by moving two
finite families to a common stage. Unique choice then turns this
bijection into an isomorphism.

Therefore, if every constructor arity \(B (a)\) is merely finite, the
polynomial functor \(S (X) = (a:A)\times (B (a)\to X)\) preserves
filtered colimits. For presentations whose construction also requires
common-stage substitutions into equations, the corresponding variable
contexts must likewise satisfy the relevant finiteness condition. The
inflationary construction can then use a simple filtered size diagram,
such as a natural-number-indexed diagram with finite upper bounds,
rather than an extensional plump ordinal basis.

These two variants establish the relevant alternatives directly. In both
cases the essential datum remains a coherent common-stage lift:
normalisation provides canonical representatives, while finiteness makes
every required family small enough for filteredness.

\section{9. Related work}\label{related-work}

Fiore, Pitts and Steenkamp introduced QW-types and their indexed
generalisation, QWI-types, as initial algebras for possibly infinitary
equational theories
\citep{fiore2020-infinitary-qiit, fiore2022-quotient-inductive}. Their
construction interleaves term formation with quotienting and obtains the
initial algebra as the colimit of a size-indexed inflationary iteration.
WISC is used to construct a plump size type sufficiently large that the
relevant polynomial functors preserve all diagrams indexed by those
sizes. Its use as a collection principle is studied in site-theoretic
and realizability settings
\citep{roberts2012-internal-categories, streicher2005-wisc}. Our
construction uses the same general inflationary framework but replaces
this uniform cocontinuity theorem with a direct proof for the particular
diagram of approximants. For depth-preserving signatures, extensional
ranks provide canonical representatives and an explicit section of the
canonical polynomial-colimit map. The result is relative to an
extensional plump ordinal basis but does not otherwise use WISC.

The restriction to a subclass of infinitary presentations is necessary
in view of known choice obstructions. Blass exhibited an infinitary
equational theory whose free algebras cannot be proved to exist in ZF
without the Axiom of Choice \citep{blass1983-free-algebra-coequalizer}.
Lumsdaine and Shulman express this theory as a higher inductive type,
showing that a sufficiently general existence principle for infinitary
higher inductive types has corresponding choice-theoretic strength
\citep[sec.~9]{lumsdaine2018-higher-inductive-types}. Our result does
not circumvent this obstruction. Instead, depth preservation identifies
structural information absent from arbitrary infinitary equational
theories and thereby isolates a class for which the necessary lifting
can be performed canonically.

Swan first introduced W-types with reductions as initial algebras for
pointed polynomial endofunctors and used them in constructions related
to the small object argument \citep{swan2018-w-types-reductions}. He
later defines \emph{image-preserving QW-types}, whose equations identify
constructor applications only when the corresponding families of
variables have the same image \citep{swan2022-hits-in-zf}. He constructs
these QW-types in classical ZF without the Axiom of Choice or an
assumption of large regular cardinals. The construction proceeds by
transfinite recursion through the class of ordinals and assigns to an
element the least ordinal stage through which the image of its
constructor arguments factors. This minimum rank is well-defined because
image-preserving equations preserve the relevant family of predecessors.
The argument uses classical properties of ordinals, including the fact
that a set of ordinals is order-isomorphic to an ordinal, and does not
directly extend to intuitionistic set theory. Our extensional plump
ordinal rank plays a role analogous to Swan's least-stage rank, but it
is formulated internally and constructively relative to the assumed
ordinal basis. Depth preservation also permits nested equations whose
two sides preserve rank without necessarily being image-preserving in
Swan's direct sense.

The original construction of infinitary QW-types by Fiore, Pitts and
Steenkamp used Abel-style sized types to express inflationary fixed
points and justify recursive definitions
\citep{abel2012-inflationary, fiore2020-infinitary-qiit}. Sized types
provide a syntactic approximation to transfinite iteration, with size
indices ensuring that recursive calls occur at smaller stages. The later
QWI development replaced this formulation with semantic size objects
constructed using WISC after logical unsoundness was discovered in the
relevant implementation of Agda sized types
\citep{fiore2022-quotient-inductive}. Our construction retains explicit
ordinal indices but makes ranks extensional so that they descend through
the quotient equations.

Several constructions apply specifically to finitary QITs and QIITs. Van
der Weide constructs finitary set-truncated higher inductive types from
quotients, using the preservation properties available for finite
arities \citep{van-der-weide2019-set-truncated-hits}. Dybjer and
Moeneclaey give a syntactic scheme for finitary higher inductive types
and interpret it in the groupoid model
\citep{dyber2018-finitary-hit-groupoid}. Kaposi, Kovács and Altenkirch
introduce a theory of signatures for finitary QIITs and construct
initial algebras and dependent eliminators for the QIITs described by
that theory, relative to the availability of the theory-of-signatures
syntax itself \citep{kaposi2019-constructing-qiit}. These finitary
results avoid the infinitary common-stage problem because finite
families can be lifted through filtered colimits using finite
directedness. Section 8.2 explains how mere finiteness of the arities
suffices constructively.

The present paper treats non-indexed QW-types. A separate line of work
studies reductions of QIITs, including internal syntax such as the
\(\mathsf{Con}\)--\(\mathsf{Ty}\) example, to simpler equational
presentations. We do not claim such an extension here, although
combining a QIIT reduction with the cocontinuity criteria developed in
this paper is a natural direction for future work.

\section{10. Conclusion}\label{conclusion}

For depth-preserving QW-signatures, an extensional plump ordinal rank
descends through the stage quotients and the colimit. Canonical
exact-rank lifting then represents every colimit element at its
intrinsic rank. Applying this lift pointwise and taking a plump supremum
constructs the section \[
\sigma_D:S \left(\mathrm{colim}_\alpha D_\alpha\right)
\longrightarrow
\mathrm{colim}_\alpha S (D_\alpha)
\] required by inflationary iteration. This yields the corresponding
initial QW-algebra from W-types and quotients, relative to the ordinal
basis and without a further use of WISC.

The construction isolates coherent sectionability, rather than depth
preservation itself, as the semantic condition. Normalising signatures
can obtain canonical lifts from normal forms, while finitary signatures
obtain common stages from finite directedness. Extending these methods
to indexed or conditional presentations and determining the exact
relationship between extensional plump ordinal bases and collection
principles remain open.

\bibliography{master.bib}

\end{document}
